\documentclass[11pt]{article}
\usepackage[margin=1in]{geometry}
\usepackage{amsmath,amssymb,amsthm}
\usepackage{booktabs}
\usepackage{graphicx}
\usepackage{hyperref}
\usepackage{microtype}
\usepackage{ctex}
\emergencystretch=2.5em

\newtheorem{definition}{定义}
\newtheorem{proposition}{命题}

\title{分形持续学习之前\\large 实现审计与预注册的决定性比较}
\author{HUSHFUSION Forge 研究}
\date{2026 年 10 月}

\begin{document}
\maketitle

\begin{quote}
\itshape
``暴力是无能者最后的手段。''
\\
—— 阿西莫夫《基地》(1951)
\end{quote}

\begin{abstract}
持续学习通常被表述为固定计算图下的参数适应; 计算图本身是否应当生长——
结构可塑性——是一个独立的问题, 常被断言而鲜被审计。我们构造了一个递归分形环学习器,
通过归一化凸门控混合根节点与后裔节点的输出, 并在改进减缓时生长后裔, 与固定单节点
学习器配对运行于重采样输入的控制环境。评审后的实现审计揭示: 在原始 72,335 轮运行中,
后裔节点从未被训练——它们不出现于任何优化器的参数表, 所宣称的深度相关稀疏调度也从未
进入任何调度判断——因此表面上的边界结果``无结构优势''并不是对分形学习器的检验。
我们修复实现、配平梯度预算, 并预注册决定性比较。本文表明: 按预注册规则, 修复后的
重跑(200,000 轮, seed 20261012)判为\emph{等价}——warm 代价差的配对 95\% bootstrap
置信区间完全落在 $\pm 1$ 步等价界内($n=5{,}000$, CI $[-0.651,-0.548]$)。
两个学习器的长程适应得到确认(warm 代价从约 130 步降至约 9 步, 对照随机基线约 32 步),
生长可操作(1 到 7 节点); 但在该确定性小模型 regime 下, 配平计算量时结构生长无可测优势。
\end{abstract}

\section{引言}

持续学习通常被表述为一个权重更新问题: 保持图结构固定, 调整参数, 使新任务被学习而
不遗忘旧任务 \cite{Finn2017,Li2017,Kirkpatrick2017}。这个框架留下了一个未被提出的问题:
当任务本身变化时, 应该改变的不仅是权重, 还有更新环路的\emph{组织}吗? 对必须在数万轮中
跟踪移动最优解的长程控制器而言, 答案至关重要; 然而结构生长是否有益, 几乎总是被断言,
而从未被测量。

我们把它当作一个测量问题。主实验让递归分形环学习器与固定单节点学习器在重采样输入的
控制环境中配对运行 72,335 轮。评审后的审计发现: 分形树的子节点从未被训练, 所宣称的
稀疏调度从未执行——长程运行检验的是门控未训练随机子网的根节点, 而非分形学习器。
修复实现并配平梯度预算(Forge commit \texttt{53e4938})后, 我们预注册判定规则并启动
200,000 轮的播种重跑。预注册判定是\emph{等价}: 配对 bootstrap 置信区间完全落在
$\pm 1$ 步界内, 因此在该 regime 下、配平计算量时, 结构生长无优势。
元学习 \cite{Finn2017,Li2017}、持续整合 \cite{Kirkpatrick2017} 与分形分配
\cite{Murray1926,LmPrincipleFractal} 的相关工作在节~\ref{sec:related} 总结;
什么成立、什么撤回、什么已判定, 是本文的完整主张集(表~\ref{tab:verdict})。

\section{结果}

\subsection{修复后的预注册重跑判定: 配平计算量下等价}

审计使结构优势处于不可判定状态(节~\ref{sec:audit}); 决定性检验需要三项与审稿人
商定、且在运行前落实的修复。第一, 计算量必须配平: 七节点树每轮的梯度步数是单节点的
数倍, 朴素重跑会把结构与计算量混淆。修复后的实现逐轮计数树中的总梯度步数, 并给固定
学习器相同预算(Forge commit \texttt{53e4938})。第二, 判定规则必须在重跑开始前定死:
等价区间 $\pm 1$ 步(度量的整数步分辨率)、收敛窗口轮 $\ge 2000$、最小 5,000 与最大
10,000 warm 轮且在 5,000 处一次窥视、以配对 bootstrap 95\% 置信区间为裁决
(细节见 \texttt{PREREGISTRATION.md})。第三, 运行必须可复现: 固定 seed(20261012)
与已提交的 trace。

重跑完成 200,000 轮(100,000 warm、100,000 cold; warm 轮 $\ge 2000$ 者
$n=99{,}000$)。在预注册的首次窥视($n=5{,}000$)处, warm 代价配对差(固定 $-$ 分形)
均值为 $-0.60$ 步, 95\% bootstrap CI 为 $[-0.651,-0.548]$——完全落在
$(-1,+1)$ 内, 即度量的分辨率内(表~\ref{tab:rerun})。按预注册规则, 此裁决为
\emph{等价}: 在一步或以上的代价分辨率下测不到结构优势。全收敛窗口确认了该窥视
($n=99{,}000$, 均值 $-0.15$), 四个时间段均值($-0.16$、$-0.12$、$-0.18$、
$-0.14$)表明差异平稳, 非早期瞬态。越界率低于预注册的 $10\%$ 上限(固定 $0.10\%$、
分形 $0.00\%$), cold 基线与预期的约 32 步一致(31.44 / 31.53), 分形树本轮同样
从 1 生长到 7 节点(封顶)。

\begin{table}[t]
\centering
\caption{预注册决定性重跑(200,000 轮, seed 20261012, 配平计算量)。warm 行为推荐器,
cold 行为随机对照。CI 为预注册首次窥视($n=5{,}000$ warm 轮, seed 固定 bootstrap)
的配对 95\% 区间。正差表示分形更好(步数更少)。裁决: \emph{等价}——CI 完全落在
$\pm 1$ 步等价界内。}
\label{tab:rerun}
\begin{tabular}{lr}
\toprule
指标 & 值 \\
\midrule
配对 warm 代价均值差 (固定 $-$ 分形), $n{=}5{,}000$ & $-0.60$ 步 \\
95\% bootstrap CI & $[-0.651,\,-0.548]$ \\
等价界 & $[-1,\,+1]$ 步 \\
收敛窗口均值差, $n{=}99{,}000$ & $-0.15$ 步 \\
warm 均值, 收敛段 (固定 / 分形) & $8.96$ / $9.11$ \\
cold 均值 (固定 / 分形) & $31.44$ / $31.53$ \\
越界率 (固定 / 分形) & $0.10\%$ / $0.00\%$ \\
最终树节点 & 7 (封顶) \\
\bottomrule
\end{tabular}
\end{table}

\subsection{审计: 后裔从未被训练}\label{sec:audit}

本文早先版本声称一个边界结果: ``无结构优势''。评审后对主实验实现的审计将该主张
撤回为\emph{这轮运行不可判定}, 因为三项发现表明这轮运行没有检验分形学习器。

\textbf{后裔从未被训练。} 在 \texttt{design\_daemon\_fractal.py}(被审快照: commit
\texttt{10ba961\textasciicircum})中, 每个节点的优化器只对其自身网络与门控参数构造
(``opt = Adam(list(self.net.parameters()) + [self.gate])'')。生长节点的子节点被追加到
\texttt{self.children}, 但从未加入任何父节点的优化器参数表, 主循环只对根调用
\texttt{train\_step}。因此后裔的参数从不接收梯度步; 每个子节点唯一的操作是 warm-up
计数器递减。所生长的``树''是一个门控固定、随机初始化子网的根网络。

\textbf{所宣称的稀疏调度从未执行。} 代码计算 \texttt{update\_freq = D**depth} 并存储,
但没有任何调度决策查询它。深度相关的更新频率是文档, 而非行为。

\textbf{评估不能按所述复现。} 72,335 轮 trace 及其分析脚本不在仓库中; daemon 以
墙钟时间为其 RNG 播种(\texttt{int(time.time()) \% 10000}); 重跑不复现报道的数字。

\textbf{后果。} 配对长程差(11.5 vs.\ 11.1 warm 步)与三 seed、701 条记录比较, 比较的
都是一个经训练的根对一个门控未训练子网的根——这一条件既不确认也不否认分形优势。
长程适应与生长检测仍然成立(节~\ref{sec:adapt}--\ref{sec:growth}); 结构优势在
节~\ref{sec:rerun} 的重跑之前一直不可判定。

\subsection{两个根学习器的长程适应均真实成立}\label{sec:adapt}

两个根都学习了并持续学习。warm 性能在前 1,000 轮从约 130 步降至大约 2,000 轮后的
约 9 步, 并在剩余 70,000 轮中保持; cold 随机基线全程保持约 32 步(表~\ref{tab:rerun})。
持续适应因此有效且稳定, 没有衰减。此结果独立于后裔训练缺陷: 它是根学习器的性质,
两个根都表现出它。

\subsection{生长可操作, 但被封顶停止}\label{sec:growth}

分形树生长了六次(第 46、57、73、80、87、94 轮), 从 1 节点到 7 节点, 然后停止。
生长由相对衰减检测器驱动(方法): 当验证损失的指数加权改进率低于该结构早期改进率的
某个比例时, 生长后裔。达到七节点后, 检测器在剩余 72,000 轮中从未再触发。我们不将
其归因于自主饱和: 七是实现中的硬安全上限, 因此生长是被封顶停止的, 而非涌现的均衡。
没有封顶时树是否会继续生长, 未被检验。审计使读法更精确: 检测器与图突变在长程上
工作, 但在被审运行中, 生长的组件从未作为学习器被发挥。

\section{讨论}

三句话必须分开, 而本文现在裁决了其中两句。``递归结构可塑性可实现且稳定''被限定:
生长事件触发, 图结构突变, 但在被审运行中生长的组件从未学习, 所以我们只能说
\emph{检测器与突变}工作, 而非生长的学习器。``结构生长改善性能''是本文早先的主张,
现朝否定方向裁决: 在配平计算量下、在该确定性控制环境中, 预注册比较判为等价——结构
生长在代价的一步分辨率下无可测优势。混淆``树生长了''与``生长的组件学习了''正是本次
审计纠正的错误; 将两者分开是本文的方法论要点。

等价裁决是边界结果, 不是普遍化的许可。它适用于小模型、确定性重采样输入环境与配平
梯度预算。分形分配论证 \cite{LmPrincipleFractal} 要求边际收益随深度急剧下降; 在数据
有限的小学习器上它们可能不会, 而在专家间分割预算可能只会稀释每位专家的证据。
分形组织在更大模型、非平稳环境或不配平预算下是否获益仍然开放——而本文的预注册协议
正是这些主张应当被检验、而非被断言的模板。

方法论的教训是程序性的。分三步评估结构生长:(i) 图是否能安全突变,(ii) 突变后学习器
是否持续适应,(iii) 突变是否在配平计算量与数据下改善预先注册的指标。本文跑了全部三步;
第 (iii) 步需要修复实现并预注册规则, 而它在该 regime 下的答案是无可测优势。

\section{相关工作}\label{sec:related}

我们的设计处于基于优化的元学习 \cite{Finn2017,Li2017}、持续学习整合
\cite{Kirkpatrick2017} 与分形分配论证的交汇点——后者根植于默里对分支网络最小功的
分析 \cite{Murray1926}, 并在我们配套的 Lean 开发中形式化 \cite{LmPrincipleFractal}。
我们将其作为动机而非预测; 本文的贡献是方法论的——一次看似干净的长程比较如何静默地
未能检验其假设, 以及决定性检验需要什么——而非声称分形组织有益或有害。

\section{方法}

\subsection{递归分形环架构}

每个节点 $v$ 持有一个预测器 $f_v(x;\theta_v)$, 并可有后裔 $\mathcal{C}(v)$。节点输出
为凸组合,
\begin{equation}
 y_v(x) = g_{v,0}\,f_v(x;\theta_v) + \sum_{u\in\mathcal{C}(v)} g_{v,u}\,y_u(x),
 \qquad g_{v,j}\ge 0,\quad \sum_j g_{v,j}=1,
\end{equation}
门控由 softmax 产生。后裔不会立即被信任: 它们在进入父节点混合之前被独立预热, 这防止
了``死子节点''失效——零初始化门控给予子节点零贡献, 因此零梯度。根节点每轮更新;
深度 $d$ 的节点以 $D^{-d}$ 的频率更新, 稀疏因子 $D$(我们用 $D=4$)。如审计所示,
该稀疏调度在部署代码中被指定但未被执行。

\subsection{生长规则}

生长由验证损失的相对衰减触发。设 $\bar\imath_t$ 为指数加权改进率, $\imath_0$ 为该结构
早期平均改进率。当 $\bar\imath_t/\imath_0$ 连续多次观测低于阈值(部署运行中为 0.25)时,
生长一个后裔; 检测器在生长后重置。干净损失过滤器排除安全惩罚值, 因为惩罚尖峰不是
结构瓶颈的证据。

\subsection{控制环境}

我们由双流环模型装配一个确定性 $\mu$-动力学控制器:
\begin{equation}
 \mu_{t+1}=\mu_t+\eta_t(1-\mu_t),
 \qquad
 \eta_t=\eta_{\mathrm{ext}}+(2\lambda-\lambda^2)(1-\mu_t),
\end{equation}
其中 $\eta_{\mathrm{ext}}$ 是外部驱动, $\lambda$ 是收敛参数。任务是达到
$\mu\ge 0.999$(工作窗口)而不越过 $\mu_{\mathrm{ceil}}=1-m_e/m_i\approx0.999780568$
(安全上限)。代价是到达窗口的步数; 越界是失败。环境是确定性的, 有清晰的安全边界,
并且是重采样的: 每轮抽取新的 $\mu_0\sim U(0.01,0.8)$。我们将其描述为重采样而非漂移:
从控制点到代价的映射是固定的, 因此学习器面对的是输入重采样而非变化的任务。它是方法
研究的基准, 而非任何实验性反引力装置的证据; $\lambda$ 在上游模型中仍是输入缺口。

\subsection{配对长程协议}

每轮两个学习器看到相同的 $\mu_0$。每个从自身经验选择控制点($(\eta_{\mathrm{ext}},\lambda)$
的 $20\times10$ 网格上的 UCB), 执行环境, 并从最近经验滑动窗(容量 400)更新参数。
奇轮为 cold 随机对照, 偶轮为 warm 推荐。部署运行使用 72,335 轮, 配失败时重启 daemon
的监督器; 决定性重跑使用 200,000 轮、seed 20261012 与配平计算量修复。如
节~\ref{sec:audit} 所述, 部署实现未将后裔参数纳入任何优化器, 因此协议描述是预期设计;
决定性重跑执行的是修复后的设计。

\subsection{配平计算量修复与预注册判定规则}

修复(Forge commit \texttt{53e4938})使每个 \texttt{train\_step} 返回其总梯度步数;
每轮, 固定学习器获得与分形树本轮相同的梯度步数, 结构与计算量不再混淆。判定规则在
重跑开始前定死(\texttt{PREREGISTRATION.md}): 收敛段(warm 轮 $\ge 2000$)的配对
warm 代价差(固定 $-$ 分形); 等价区间 $\pm 1$ 步; 配对 bootstrap 95\% CI($10{,}000$
次重采样, seed 固定); 裁决为等价 / 分形更优 / 固定更优 / 未决, 在 $n=5{,}000$ 处
一次窥视, 若仍未决则在 $n=10{,}000$ 处停止, 不再追加样本(防止持续窥视)。

\subsection{塑造设计的失败分析}

四个失败值得记录。(i) 绝对``无改进''生长检测器在浅层网络上从不触发(它们缓慢地永远
改进); 相对衰减检测器修复了它。(ii) 零初始化的子门控产生死子节点; 均匀门控初始化加
预热阶段修复了它。(iii) 生长一层而不重新初始化逐参数步长元参数, 损害了快速适应环,
因此结构生长与元参数生长不能视为独立操作。(iv) 标准 RSI 指标不适用于单调改进的损失
轨迹(单侧增量使 RSI $\equiv100$), 因此此处该词用于``递归自我改进'', 而非金融指标。

\section*{数据与代码可用性}
被审实现、修复与分析脚本在 HUSHFUSION Forge 仓库中 \cite{HushfusionForge}。被审快照为
commit \texttt{10ba961\textasciicircum}(子节点未训练); 修复后的配平计算量版本为
commit \texttt{53e4938}。决定性重跑 trace 为 \texttt{oml\_daemon\_fractal\_matched\_20261012.jsonl}
(200,000 轮, seed 20261012; 36,151,557 字节; md5
\texttt{edc0f8e78fb0527d29d151222473f26b}; 最后写入 2026-10-12T03:11 UTC), 由
\texttt{paper/analyze\_matched.py} 分析; 预注册、判定规则与裁决见 \texttt{PREREGISTRATION.md}。
写作标准见 \texttt{WRITING-STANDARD.md}。

\bibliographystyle{plain}
\bibliography{main}

\end{document}